\section{Complete calibration of the private-binder transfer source}
\label{sec:genc-calibration}

The source \file{source/gc215.sites.json} uses the original single
cache guard.  This section proves its complete command-level calibration.
Its only logical-interface change is that the ninth
AX output generalizes with the independently supplied private parameter
$C$, rather than the incoming counter $U$.  The argument below supplies
both directions of the refutation equivalence, the exact wire packing,
and the protected-tail bridge needed for the measured blank initialization.
The source and binary runtime contracts are separate, explicitly identified
premises; finite execution tests are not used to infer these unbounded
mathematical statements.

\subsection{The retained numerical calculus}

For all natural numbers define
\[
 P(a,b)=\frac{(a+b)(a+b+1)}2+b,
 \quad I(a,b)=4P(a,b),\quad a\to b=I(a,b),
\]
and
\[
 R(a,b)=4P(a,b)+1,\qquad Q_a b=4P(a,b)+2,
 \qquad E(a,b)=4P(a,b)+3.
\]
Truth is $0$, the refutation target is $1$, $NX=I(X,1)$, and
$\exists_aX=N(Q_a(NX))$.  In particular $I(0,0)=0$ and $R(0,0)=1$.
Every parameter ranges over every natural number, including numbers that
are not conventional syntactic formula codes.

Let $\mathcal D$ be the inherited dominated native calculus.  Its rules
are truth, exact numerical modus ponens, and theorem-only generalization:
\[
 \mathcal D\vdash0,\qquad
 \frac{\mathcal D\vdash X\quad\mathcal D\vdash I(X,Y)}
      {\mathcal D\vdash Y},\qquad
 \frac{\mathcal D\vdash X}{\mathcal D\vdash Q_iX}.
\]
For specificity, write the seven \emph{emitted} row functions as
\begin{align*}
 R_1(a,b,c)&=\exists_aE(a,b),\\
 R_2(a,b,c)&=Q_a(b\to c)\to(Q_ab\to Q_ac),\\
 R_3(a,b,c)&=E(a,b)\to(R(c,a)\to R(c,b)),\\
 R_4(a,b,c)&=Q_aQ_bc\to Q_bQ_ac,\\
 R_5(a,b,c)&=E(c,a)\to(R(c,b)\to R(a,b)),\\
 R_6(a,b,c)&=(\exists_aQ_ac)\to c,\\
 R_7(a,b,c)&=E(c,a)\to\bigl(E(c,b)\to Q_{R(a,b)}E(a,b)\bigr).
\end{align*}
The middle argument of $R_6$ is unused.  This explicitly respects the
repository's correction that its sixth output is $\rho_6(a,c)$.
The eighth family is
\[
 L(A,B,C,U)=((C\to U)\to A)\to((A\to C)\to(B\to C)),
 \qquad U\le A.
\]
The full inherited gated seed bundle is retained literally.  To specify it,
put
\begin{align*}
 \mathcal B(b,k,m)&=Q_b\bigl(m\to NQ_b(m\to NR(b,k))\bigr),\\
 \mathcal S(a,b,k)&=Q_k(R(k,a)\to R(k,b)),\\
 \mathcal C(a,b,k,m)&=Q_b\bigl(R(b,k)\to
      NQ_a(m\to NQ_a(\mathcal S(a,b,k)\to\mathcal B(b,k,m)))\bigr),\\
 \mathcal J(a,b,c)&=Q_c\bigl(\mathcal C(a,b,c,c)
                              \to N\mathcal B(b,c,c)\bigr),\\
 \mathcal V(a,b,c)&=R(a,b)\to Q_cR(a,b),\\
 \operatorname{Bundle}(a,b,c)&=N\bigl(\mathcal V(a,b,c)
                                          \to\mathcal J(a,b,c)\bigr).
\end{align*}
It is an axiom when $c\ne a$ and $c\ne b$.  Define the gated output
$\Gamma(a,b,c)$ to be this bundle when both inequalities hold and $0$
otherwise.  No support-bearing part of this bundle is deleted.

The previously established calibration of $\mathcal D$ is used as a
premise~\cite{dominated,st215}: ZFC proves that a finite $\mathcal D$-derivation of $1$ exists
if and only if a finite ZFC refutation exists.  This premise includes the
inherited all-natural certificate recovering unrestricted $L$ from its
dominated instances.  The present change does not alter any axiom or rule
of $\mathcal D$; it changes how one command samples theorem generalization.

\subsection{Exact commands and field order}

At a command boundary use $(t,w,d,h)$ for top, numerical stack, additive
counter, and cache.  Write $[x;w]=P(x,w)$ and use right nesting for a
longer explicit list.  For every natural $w$, let
$(\operatorname{head}(w),\operatorname{tail}(w))$ be its Cantor unpairing;
in particular the two projections of zero are zero.

The exact transfer command is
\[
 S(t,w,d,h)=
 \begin{cases}
  \mathsf{halt},&h=1,\\
  (h,\operatorname{tail}(w),d,\operatorname{head}(w)),&h\ne1.
 \end{cases}
\]
It tests the \emph{incoming cache}, not the incoming top or the newly
popped head.  For $v=d+x$, the exact MP command is
\[
 \mathsf{MP}(x)(t,w,d,h)=
 \begin{cases}
  (v,[0;w],v,0),&I(t,v)=h,\\
  (0,w,v,0),&I(t,v)\ne h.
 \end{cases}
\]
For $\mathsf{AX}(c;x,y)$ put
\[
 U=d,\qquad A=d+x,\qquad B=h+y,\qquad C=c.
\]
The new exact AX equation is
\begin{align*}
 t'&=\Gamma(A,B,C),&d'&=0,&h'&=0,\\
 w'&=[R_1(A,B,C),\ldots,R_7(A,B,C),
          L(A,B,C,U),Q_C h,t;w].
\end{align*}
The $h$ inside $Q_Ch$ is the incoming cache, before adding $y$.
Only this ninth head differs from the parent interface, whose ninth head
was $Q_Uh$.  The inequality $U\le A$ holds on every natural input.
The gate and first eight heads are implication codes and hence cannot
equal $1$.  The ninth head has residue $2$ modulo $4$, so it is neither
$0$ nor $1$.  The saved incoming top is the tenth head.

The actual field order is
\[
 \operatorname{fields}(S)=[0],\qquad
 \operatorname{fields}(\mathsf{MP}(x))=[1,x],\qquad
 \operatorname{fields}(\mathsf{AX}(c;x,y))=[c+2,y,x].
\]
The notation after the semicolon lists the counter increment before the
cache increment; the wire reads the cache increment first.  These are the
old YX fields, not an XY variant.

\subsection{Structural macros and a strong proof compiler}

Define
\[
 D=\mathsf{AX}(0;0,0),\qquad H=D;S^{11},\qquad
 \mathsf{PUSH}=D;S^9;\mathsf{MP}(0).
\]
Concatenation denotes execution in the displayed order.  For every natural
input these macros terminate, possibly by the stated halt, and satisfy
\[
 H(t,w,d,h)=
 \begin{cases}
  \mathsf{halt},&t=1,\\
  (t,\operatorname{tail}(w),0,\operatorname{head}(w)),&t\ne1,
 \end{cases}
 \qquad
 \mathsf{PUSH}(t,w,d,h)=(0,[t;w],0,0).
\]
To verify both equations, let $a_1,\ldots,a_8,q,t$ be the ten AX
heads, where now $q=Q_0h$.  The first ten shifts test only $0$,
$a_1,\ldots,a_8,q$, all different from $1$.  After ten shifts the
tuple is $(q,w,0,t)$.  The eleventh therefore has exactly the displayed
top-falsity guard and pop effect.  After nine shifts, instead, the tuple
is $(a_8,[t;w],0,q)$.  Its final MP necessarily fails because the cache
has tag two and every $I(a_8,0)$ has tag zero.  Thus PUSH never halts,
even when the original top or cache is $1$.  In particular, $H$ is not
silently identified with the parent's cache-guarded NIP.

A strong block $B_X$ will satisfy, for \emph{every} pair of natural
numbers $t,w$,
\[
 B_X(t,w,0,0)=\mathsf{halt}\quad\hbox{or}\quad(X,w,0,0).
 \tag{SC}
\]
No theoremhood assumption is made about this arbitrary input top or stack.
Define
\[
 \mathsf{DISCARD}=\mathsf{MP}(1);\mathsf{AX}(1;0,0);
                      S^{10};\mathsf{MP}(0).
\]
On $(t,w,0,0)$ the first MP fails, since $I(t,1)>0$, and returns
$(0,w,1,0)$.  The AX has $A=C=1$, so its gate is zero.  Ten shifts
give $(Q_1 0,w,0,0)$.  Since $Q_1 0>0$, the final MP fails and returns
$(0,w,0,0)$.  This proves (SC) for truth, with no halt.
On a cleared tuple $(0,w,0,0)$ put
\[
 J(u)=
 \begin{cases}
  \text{empty word},&u=0,\\
  \mathsf{MP}(u),&u>0.
 \end{cases}
\]
It returns $(0,w,u,0)$.  The empty case is essential: $\mathsf{MP}(0)$
would succeed by $I(0,0)=0$ and introduce an unwanted explicit zero head.

\paragraph{Axiom retrieval.}
For a retained row at position $1\le j\le7$, use
\[
 B_{R_j(a,b,c)}=
 \mathsf{DISCARD};\mathsf{AX}(c;a,b);S^{j+1};H^{9-j}.
\]
After $j+1$ shifts the selected row is top, and exactly $9-j$ of
the ten generated heads remain above the original stack.  The last is
the saved zero.  Each $H$ preserves the selected, non-$1$ implication
code and removes one of these heads.  The final cache is therefore zero
and the original stack is restored exactly.  For any desired
$L(a,b,c,u)$ with $u\le a$, use
\[
 B_{L(a,b,c,u)}=
 \mathsf{DISCARD};J(u);\mathsf{AX}(c;a-u,b);S^9;H.
\]
Here $U=u$, $A=a$, $B=b$, and $C=c$; the single $H$ removes saved zero.
For an enabled full bundle use
\[
 B_{\operatorname{Bundle}(a,b,c)}=
 \mathsf{DISCARD};\mathsf{AX}(c;a,b);H^{10}.
\]
Its implication-coded top is preserved as all ten generated heads are
removed.  These calculations apply to exact native codes, including
zero-valued row instances.

\paragraph{The new generalization block.}
Given $B_X$, set
\[
 B_{Q_iX}=B_X;\mathsf{PUSH};S;\mathsf{AX}(i;0,0);S^{10}.
 \tag{GEN-C}
\]
If the premise block returns, PUSH gives $(0,[X;w],0,0)$ and the single
shift gives $(0,w,0,X)$.  That shift tests incoming cache zero, so it is
valid also for $X=1$.  AX now has $C=i$ and its ninth head is exactly
$Q_iX$, followed by saved zero.  Ten shifts return $(Q_iX,w,0,0)$.
This is precisely where using the private binder changes the compiler:
the parent's $J(i)$ is omitted.  Here $A=0$ and $B=X$, so the AX gate
need not vanish when $i$ differs from both.  It is nevertheless an
implication code, and the first shift safely overwrites it.  The proof
does not rely on a false zero-gate assertion.

\paragraph{Modus ponens, with the major compiled first.}
Let $M=I(A,B)$, and assume strong blocks for $A$ and $M$.  Put
\[
 B_B=B_M;\mathsf{PUSH};B_A;\mathsf{PUSH};S;S;
       \mathsf{MP}(B);\mathsf{AX}(B;0,0);S^{11}.
 \tag{MP-C}
\]
Unless an earlier permitted halt occurs, the first four blocks give
$(0,[A,M;w],0,0)$.  Two shifts either halt on $A=1$ or give
$(A,w,0,M)$.  The next MP succeeds by literal numerical equality and
returns $(B,[0;w],B,0)$.  The cleanup AX has $A=C=B$, hence zero
gate, and its saved tenth head is $B$.  After ten shifts the tuple is
$(Q_B0,[0;w],0,B)$.  The eleventh halts if $B=1$, or returns
$(B,w,0,0)$ otherwise.  This proves (SC).  When $A=B=0$, also $M=0$;
the successful zero-alias MP still intentionally pushes one explicit
zero, and the cleanup removes exactly that zero.  No nonzero-formula
assumption is hidden in the argument.

Induction on a finite native derivation now gives (SC) for every theorem.
A finite proof DAG may first be expanded into a finite proof tree;
compiling the major before the minor introduces no circular dependency.
If $\mathcal D\vdash1$, the fixed finite word
\[
                         B_1;H
\]
halts from every $(t,w,0,0)$: either $B_1$ already halts, or its result
has top $1$, which the eleventh transfer of $H$ detects.

\subsection{Exact packing and cofinal successful candidates}

For a finite field list $f$ and residual $r$, let $F(f,r)$ be its
right-nested Cantor pairing above $r$.  Define the encoded proof tail by
\[
 \operatorname{tailcode}([])=0,\qquad
 \operatorname{tailcode}(k::L)
   =1+F\bigl(\operatorname{fields}(k),\operatorname{tailcode}(L)\bigr).
\]
At a positive proof pointer the source first decrements, then reads the
command's one, two, or three fields.  Thus a positive encoded tail
executes its next command and leaves the indicated successor tail.
A freshly fetched candidate has no outer successor: its initial command
is decoded directly.

Let $r=\operatorname{tailcode}(B_1;H)$.  For every natural $Y$ choose
\[
                 N_Y=P\bigl(2,P(Y,P(1,r))\bigr).
 \tag{Cofinal}
\]
With the actual YX field order this first executes
$\mathsf{AX}(0;1,Y)$ and leaves residual $r$.  The first AX terminates
and resets counter and cache to zero.  Its resulting top and stack may
be arbitrary; the universal contract of $B_1;H$ still ensures a halt.
Cantor domination gives $N_Y\ge Y$, so successful candidates lie beyond
every finite startup cutoff.

Every nonhalting positive proof pointer $p$ drains strictly: after its
initial decrement, the next residual is obtained by right projections
of $p-1$, so it is at most $p-1<p$.  Each source command terminates or
halts by its all-natural macro contract.  Therefore a candidate cannot
block the scheduler forever.  The next-candidate counter increases by
one, so the scheduler eventually reaches any chosen $N_Y$ above its
initial cutoff, unless it halts earlier.  This proves completeness
without a finite search bound or an additional fairness assumption.

\subsection{An all-natural protected-tail theorem}

The actual blank initializer enters a partial AX computation, not an
ordinary zero-stack root.  The following schedule theorem proves that
the arbitrary numerical suffix underneath its safe heads is never exposed.
It depends only on the exact three field lengths and command balances,
so changing the generalization binder does not change its combinatorics.

Let $f$ be the right Cantor projection and $T_s=s(s+1)/2$.  Let $W(n)$
be the finite syntactic command word parsed from fresh candidate $n$,
ignoring possible semantic halts, and put $\mathcal R(0)=[]$ and
$\mathcal R(r)=W(r-1)$ for $r>0$.  If
$\operatorname{unpair}(n)=(a,y)$, then exactly
\[
 W(n)=
 \begin{cases}
  S;\mathcal R(y),&a=0,\\
  \mathsf{MP};\mathcal R(f(y)),&a=1,\\
  \mathsf{AX};\mathcal R(f(f(y))),&a\ge2.
 \end{cases}
 \tag{Schedule}
\]
Payloads are still read; they are irrelevant only to the following
conservative balance.  Give $S,\mathsf{MP},\mathsf{AX}$ balances
$-1,0,+10$, respectively.  Successful MP actually adds a further
theorem zero, which this estimate may safely ignore.  Write $b(n)$
for the total balance of $W(n)$ and
$B(k)=\sum_{n<k}b(n)$.  Define $g(0)=0$ and $g(r)=b(r-1)$ for $r>0$.
Strict residual descent makes these definitions well founded.  Induction
also proves that $W(n)$ has at most $n+1$ commands: its residual is at
most $n$.  Consequently every prefix of $\mathcal R(r)$ has balance
at least $-r$.

\paragraph{Prefix pullback.}
If a sequence $q(0),q(1),\ldots$ has nonnegative initial partial sums,
then the sequence $q(f(0)),q(f(1)),\ldots$ does too.  Indeed, increasing
inputs to $f$ enumerate the consecutive blocks
$[0],[0,1],[0,1,2],\ldots$.  Any initial segment consists of complete
such blocks and an initial part of the next one.  Apply this observation
twice for $q(f(f(y)))$.  We use $q(r)=g(r)+1$, for which
\[
 \sum_{r<k}q(r)=1+B(k-1)+(k-1)\qquad(k\ge1).
 \tag{Shift}
\]
Thus previously established inequalities $B(j)+j\ge0$ supply the
needed nonnegative partial sums.  Using unshifted $g$ here would be
invalid, since the new tag order has negative initial balances.

\paragraph{Every complete diagonal from $s=2$ is prefix nonnegative.}
We prove this jointly with $B(k)+k\ge0$ and $B(T_s)\ge-3$, by induction
on the Cantor diagonal.  The base is
\[
 W(0)=[S],\qquad W(1)=[\mathsf{MP}],\qquad W(2)=[S,S],
\]
giving $B(0),B(1),B(2),B(3)=0,-1,-1,-3$.
For $s\ge2$, the candidates $T_s+y$ have selector $s-y$.
The first $s-1$ start with AX, the next with MP, and the last with S.
The first $k$ AX candidates have total balance
\[
 10k+\sum_{y<k}g(f(f(y)))\ge9k,\qquad 0\le k\le s-1,
\]
by prefix pullback and (Shift).  All required earlier indices are at
most $s-2<T_s$, so this does not assume the induction conclusion.
Before AX candidate $y$ the balance is at least $9y$; its AX adds $10$
and its suffix loses at most $f(f(y))\le y$ at any prefix.  Its
nonempty prefixes therefore have balance at least $8y+10>0$.
After the AX portion the balance is at least $9(s-1)$.  The MP
candidate's suffix loses at most $s-1$, leaving at least $8(s-1)$
throughout.  The last S costs one and its suffix loses at most $s$,
leaving at least
\[
                         8(s-1)-1-s=7s-9>0.
\]
Thus the entire diagonal is prefix nonnegative and has nonnegative final
balance.  Since its starting $B$ is at least $-3$, all its candidate
boundaries have $B(k)\ge-3$; their $k\ge3$ then gives $B(k)+k\ge0$.
This closes the simultaneous induction for every natural diagonal and
every command prefix, including recursively decoded suffixes.

\paragraph{Exact conservative credit at any cutoff.}
Given $N$, choose $T_s\le N<T_{s+1}$ and set
$E(N)=\max(3,T_{s+1})$.  Let $d(N)$ be the negative of the minimum
prefix balance, including the empty prefix, of the finite word
\[
                    W(N);W(N+1);\cdots;W(E(N)-1).
\]
Every later complete diagonal is prefix nonnegative.  Hence $d(N)$
is exactly the conservative credit needed for the whole infinite
syntactic schedule beginning at $N$, not an estimate extrapolated from
a finite enumeration.  With an initial positive proof pointer $p$,
let $b,m$ be the total and minimum-prefix balances of $\mathcal R(p)$.
The required credit is
\[
                         \max(0,-m,d(N)-b).
\]
In particular $E(1)=3$ and $W(1);W(2)=[\mathsf{MP},S,S]$, so
\[
                              d(1)=2.
 \tag{Credit}
\]

\paragraph{Protection and soundness.}
Represent the stack as a finite explicit list of native theorem heads
above a fixed arbitrary numerical tail $J$.  Suppose top and cache are
native theorems and the explicit list has at least the required credit.
Returning AX adds ten theorem heads: eight axiom instances, $Q_Ch$ by
theorem generalization, and the saved theorem top.  Returning MP either
preserves the list or adds theorem zero.  Before each returning S the
conservative explicit length is at least one, since its subsequent
length is nonnegative.  Hence it removes an explicit theorem head and
never decodes $J$.  It transfers theorem codes; if it halts, its incoming
cache $1$ is already a theorem.  This proves that every halt on the
protected run supplies a native refutation.

The explicit length is a proof invariant, not the size of a numerical
list encoding.  Although $P(0,0)=0$, the identity
$\operatorname{unpair}(P(x,v))=(x,v)$ holds for all $x,v$.  An explicit
zero wrapper may therefore be counted, pushed, and popped even when its
integer representation coincides with the underlying zero tail.  No
injective encoding of finite list length is assumed.  Defining $W$ while
ignoring halts also introduces no nonhalting assumption: when a halt
occurs, the argument applies to its finite executed prefix.

\paragraph{Eight safe heads suffice after an arbitrary AX.}
There is a useful stronger bootstrap variant.  Even if the incoming
top and cache of the \emph{first} AX are arbitrary, its gate is an axiom
or truth, its new cache is zero, and its first eight heads are valid
native axioms.  The ninth head $Q_C(\mathrm{old}\ h)$ need not be a
theorem.  Absorb it, the saved old top, and the whole old stack into
the single arbitrary tail
\[
                     J=[Q_C(\mathrm{old}\ h),\mathrm{old}\ t;
                                                       \mathrm{old}\ w].
\]
There are then eight explicit safe heads above $J$.  At $N=1,p=0$,
(Credit) protects that suffix with a permanent conservative margin of
at least six.  One must not instead assume that generalizing an arbitrary
dirty cache makes it a theorem.

\subsection{Application to the exact measured blank frame}

The separately replayed bootstrap receipt for this candidate gives the
following complete correspondence.  After $68{,}710$ binary transitions
from a blank tape, the machine is at dispatcher state $84$, head $-2$,
header $62$, gap zero, and source address
\texttt{00110010010110010}, which is primitive ordinal $58$.  Its logical
registers, in the recorded order $(N,T,W,A,D,H,X,Y,L)$, are
\[
                         (1,0,0,0,0,0,0,18,0),
\]
with all $63$ spare registers zero.  The physical register named $A$
in this tuple carries the private parameter called $C$ in the AX
equations; it is distinct from their derived formula parameter $A=d+x$.
The receipt compares the entire
finite set of $100$ one-cells, not merely those displayed register values.
Its independently replayed source trace starts at AX entry ordinal $25$
with
\[
                         (1,0,2,0,0,0,0,0,0)
\]
and reaches the same source ordinal and registers after exactly $177$
source primitives.  These concrete execution facts are premises supplied
by the bound trace and source-CFG checks, not conclusions of the present
logical argument.

This synthetic entry is $\mathsf{AX}(0;0,0)$ at work tuple
$(t,w,d,h)=(0,2,0,0)$, with next candidate $N=1$ and empty proof pointer.
By the all-natural AX contract, its remaining computation terminates and
reaches
\[
 t=d=h=0,\qquad
 w=[R_1(0,0,0),\ldots,R_7(0,0,0),L(0,0,0,0),Q_00,0;2].
\]
All ten displayed heads are native theorems.  The gate is zero because
$C=A=B=0$.  The ninth head is $Q_00$ in both parent and new interfaces,
but the generalization proof above covers every other input as well.
The numerical suffix $2$ must remain protected: its repeated unpairing
eventually yields falsity, since $2=P(0,1)$ and $1=P(1,0)$.
The next-candidate value is still $1$ and the proof pointer is still zero.
The required credit is only two, so the ten safe heads leave a permanent
conservative margin of at least eight.  Alternatively the general
eight-head argument applies with margin six without using theoremhood of
the saved input fields.

Consequently the partial-AX frame suffices: the all-natural source/runtime
contracts advance it to the protected root without literally evaluating
its enormous later formula codes.  A coincident register projection alone
would not justify this step; the complete frame, source ordinal, scheduler
fields, spare registers, and whole-tape trace are the separately checked
physical premises.

\paragraph{Calibration theorem and formalization scope.}
Assume the exact source command equations, the all-natural binary
simulation of that source, and the complete blank-to-source frame above.
Then ZFC proves
\[
 \text{the new blank-tape machine halts}
 \quad\Longleftrightarrow\quad
 \mathcal D\vdash1
 \quad\Longleftrightarrow\quad
 \neg\operatorname{Con}(\mathrm{ZFC}).
\]
The forward implication is the protected-tail theorem; the reverse is
the strong compiler and cofinal candidate construction.  Both are uniform
finite arithmetic arguments and combine with the inherited native
calibration inside ZFC.  Taking negations gives the requested equivalence
between nonhalting and $\operatorname{Con}(\mathrm{ZFC})$.
This is an ordinary mathematical proof tied to the new literal interface
and its separately checked physical certificates.  It does not assert a
new Lean or Rocq kernel endpoint for this exact transition table merely
because an endpoint exists for an earlier 215-state table.
